% !TEX program = xelatex
% =====================================================================
%  BÁO CÁO ĐỒ ÁN TRÍ TUỆ NHÂN TẠO, HỌC PHẦN AIP202
%  Khoa Công nghệ Thông tin, Trường Đại học Kiến trúc Đà Nẵng
%
%  BẢN HOÀN CHỈNH ĐỘC LẬP (STANDALONE) - BIÊN DỊCH TRỰC TIẾP TRÊN OVERLEAF
%  HƯỚNG DẪN: Menu > Compiler > XeLaTeX, rồi bấm Recompile (Ctrl + Enter)
% =====================================================================
\documentclass[12pt, a4paper, oneside]{book}
\usepackage[margin=2.2cm, top=2.5cm, bottom=2.5cm]{geometry}

% --- CẤU HÌNH PHÔNG CHỮ & ĐA NGÔN NGỮ (TƯƠNG THÍCH 100% OVERLEAF LINUX) ---
\usepackage{fontspec}
\usepackage[vietnamese, english, bidi=basic, provide=*]{babel}
\babelprovide[import, onchar=ids fonts]{vietnamese}
\babelprovide[import, onchar=ids fonts]{english}
\babelfont{rm}{Noto Serif}
\babelfont{sf}{Noto Sans}
\babelfont{tt}{Noto Sans Mono}

% --- CÁC GÓI TOÁN HỌC, ĐỒ HOẠ, BẢNG BIỂU & ĐỊNH DẠNG ---
\usepackage{amsmath, amssymb}
\usepackage{titlesec}
\usepackage{array}
\usepackage{longtable}
\usepackage{tabularx}
\usepackage{enumitem}
\usepackage[table]{xcolor}
\usepackage{booktabs}
\usepackage{listings}
\usepackage{float}
\usepackage{caption}
\usepackage{adjustbox}
\usepackage[skip=6pt, indent=0pt]{parskip}
\usepackage{fancyhdr}
\usepackage{hyperref}

\hypersetup{
    colorlinks=true,
    linkcolor=black,
    urlcolor=blue,
    citecolor=blue,
    pdftitle={Báo Cáo Đồ Án Trí Tuệ Nhân Tạo - Restaurant Intelligence},
    pdfauthor={Nguyễn Ngọc Thắng - Trần Lê Quốc Thịnh}
}

% --- CẤU HÌNH HEADER & FOOTER VỚI FANCYHDR ---
\newcommand{\dauPageStyle}{%
    \pagestyle{fancy}%
    \fancyhf{}%
    \fancyhead[L]{\small\itshape AIP202 -- Hệ thống Restaurant Intelligence}%
    \fancyhead[R]{\small\itshape Khoa CNTT -- ĐH Kiến Trúc Đà Nẵng}%
    \fancyfoot[C]{\thepage}%
    \renewcommand{\headrulewidth}{0.4pt}%
    \renewcommand{\footrulewidth}{0.4pt}%
}

% Cấu hình hiển thị mã nguồn / Prompt
\lstset{
    basicstyle=\ttfamily\footnotesize,
    backgroundcolor=\color{gray!8},
    frame=single,
    rulecolor=\color{gray!40},
    breaklines=true,
    postbreak=\mbox{\textcolor{red}{$\hookrightarrow$}\space},
    tabsize=2,
    showstringspaces=false,
    numbers=left,
    numberstyle=\tiny\color{gray},
    xleftmargin=1.5em,
    framexleftmargin=1.5em
}

% --- THƯ VIỆN VẼ SƠ ĐỒ (TIKZ) ---
\usepackage{tikz}
\usetikzlibrary{shapes.geometric, shapes.multipart, arrows.meta, positioning, calc, fit, backgrounds}

% Các lệnh metadata hỗ trợ tương thích mẫu của trường
\newcommand{\crname}[1]{}
\newcommand{\cmonhoc}[1]{}
\newcommand{\ctname}[1]{}
\newcommand{\cstuname}[1]{}
\newcommand{\cmssv}[1]{}
\newcommand{\clopsv}[1]{}
\newcommand{\csCouncil}[1]{}
\newcommand{\csSupervise}[1]{}
\newcommand{\cttime}[1]{}
\newcommand{\csdeptname}[1]{}
\newcommand{\csplace}[1]{}

% Định nghĩa trang bìa trang trọng có khung kép chuẩn ĐH Kiến Trúc Đà Nẵng
\newcommand{\coverpage}{%
\begin{titlepage}
\begin{tikzpicture}[remember picture, overlay]
    \draw[line width=2.5pt, draw=blue!85!black] 
        ($(current page.north west) + (1.8cm, -1.8cm)$) 
        rectangle 
        ($(current page.south east) + (-1.5cm, 1.8cm)$);
    \draw[line width=0.8pt, draw=blue!65!black] 
        ($(current page.north west) + (1.95cm, -1.95cm)$) 
        rectangle 
        ($(current page.south east) + (-1.65cm, 1.95cm)$);
\end{tikzpicture}

\begin{center}
\vspace*{0.4cm}
{\large\bfseries TRƯỜNG ĐẠI HỌC KIẾN TRÚC ĐÀ NẴNG}\\[0.2cm]
{\large\bfseries KHOA CÔNG NGHỆ THÔNG TIN}\\[0.3cm]
\rule{0.48\textwidth}{1.0pt}

\vspace{3.2cm}

{\LARGE\bfseries BÁO CÁO ĐỒ ÁN}\\[0.4cm]
{\Large\bfseries MÔN: ĐỒ ÁN TRÍ TUỆ NHÂN TẠO}\\[0.2cm]
{\normalsize\bfseries MÃ HỌC PHẦN: AIP202 \textperiodcentered{} 2 TÍN CHỈ}

\vspace{1.5cm}

{\Large\bfseries ĐỀ TÀI:\\[0.35cm]
\scalebox{0.95}{\parbox{\linewidth}{\centering HỆ THỐNG PHÂN TÍCH ĐÁNH GIÁ NHÀ HÀNG\\ỨNG DỤNG TRÍ TUỆ NHÂN TẠO\\(RESTAURANT INTELLIGENCE)}}}

\vspace{2.5cm}

\begin{tabular}{ll}
\textbf{Giảng viên hướng dẫn:} & \textbf{TS. Đỗ Phúc Hảo} \\[12pt]
\textbf{Sinh viên thực hiện:}  & \textbf{Nguyễn Ngọc Thắng} -- MSSV: 2351220287 \\[8pt]
                               & \textbf{Trần Lê Quốc Thịnh} -- MSSV: 2351220257 \\[12pt]
\textbf{Lớp / Ngành:}          & \textbf{CNTT 23 -- Công nghệ thông tin} \\
\end{tabular}

\vfill

{\bfseries Đà Nẵng, năm 2026}
\vspace*{0.3cm}
\end{center}
\end{titlepage}
}

\begin{document}

% Trang bìa trang trọng có khung kép theo quy cách trường
\coverpage

% =====================================================================
%  PHẦN ĐẦU: ĐÁNH SỐ TRANG LA MÃ THƯỜNG (i, ii, iii...)
% =====================================================================
\frontmatter
\dauPageStyle

% --- LỜI CAM ĐOAN ---
\chapter*{Lời cam đoan}
\addcontentsline{toc}{chapter}{Lời cam đoan}

Chúng em xin cam đoan đề tài nghiên cứu và xây dựng \textbf{``Hệ thống Phân tích Đánh giá Nhà hàng ứng dụng Trí tuệ nhân tạo (Restaurant Intelligence)''} là công trình khoa học do nhóm tự nghiên cứu và phát triển dưới sự hướng dẫn chuyên môn của thầy \textbf{TS. Đỗ Phúc Hảo}.

Toàn bộ số liệu thực nghiệm trên 50 nhà hàng chất lượng cao, 677 bài đánh giá thực tế đã qua tiền xử lý làm sạch, 1.842 nhận định khía cạnh trích xuất bằng mô hình ngôn ngữ lớn và kết quả 173 kịch bản kiểm thử tự động đạt 100\% PASS được ghi nhận trung thực từ hệ thống vận hành thực tế. Mọi tài liệu tham khảo, nền tảng mã nguồn mở và dữ liệu cộng đồng đều được nhóm trích dẫn minh bạch theo đúng quy chuẩn học thuật.

\vspace{1.2cm}
\begin{flushright}
\textit{Đà Nẵng, ngày 27 tháng 10 năm 2026}\\
\textbf{Nhóm sinh viên thực hiện}\\[1.2cm]
\textbf{Nguyễn Ngọc Thắng -- Trần Lê Quốc Thịnh}
\end{flushright}

% --- LỜI CẢM ƠN ---
\chapter*{Lời cảm ơn}
\addcontentsline{toc}{chapter}{Lời cảm ơn}

Để hoàn thành đề tài đồ án này, chúng em xin gửi lời cảm ơn chân thành và sâu sắc nhất đến các thầy cô giáo Khoa Công nghệ Thông tin – Trường Đại học Kiến Trúc Đà Nẵng đã trang bị cho chúng em nền tảng kiến thức công nghệ vững vàng trong suốt quá trình học tập.

Đặc biệt, chúng em xin bày tỏ lòng biết ơn sâu sắc đến thầy \textbf{TS. Đỗ Phúc Hảo} – giảng viên hướng dẫn môn học Đồ án Trí tuệ nhân tạo (AIP202). Thầy đã luôn định hướng, tận tình chỉ dẫn và đưa ra những nhận xét chuyên môn quý báu giúp nhóm hoàn thiện đề tài từ khâu thiết kế kiến trúc, xử lý dữ liệu tiếng Việt đến triển khai mô hình ngôn ngữ lớn và gia cố bảo mật.

\vspace{1.2cm}
\begin{flushright}
\textit{Đà Nẵng, năm 2026}\\
\textbf{Nhóm sinh viên thực hiện}
\end{flushright}

% --- TÓM TẮT ĐỒ ÁN ---
\chapter*{Tóm tắt đồ án}
\addcontentsline{toc}{chapter}{Tóm tắt đồ án}

Đề tài \textbf{``Restaurant Intelligence''} giải quyết triệt để bài toán phân mảnh dữ liệu ẩm thực giữa Foody.vn và Google Maps, xử lý vấn nạn đánh giá rác tại Việt Nam và bóc tách cảm xúc đa khía cạnh (Aspect-Based Sentiment Analysis -- ABSA).

Hệ thống tích hợp quy trình xử lý thông minh khép kín: thu thập dữ liệu đa nguồn bằng Playwright và Selenium; tiền xử lý lọc spam đánh giá 11 trường hợp tiết kiệm 18.6\% token API; thuật toán khớp nối thực thể ($\ge 85\%$) đối chiếu địa chỉ hành chính và mã băm SHA-256 Fingerprint chống trùng lặp; ứng dụng mô hình ngôn ngữ lớn Google Gemini bóc tách 5 khía cạnh ẩm thực cốt lõi kèm cờ cảnh báo khẩn cấp \texttt{is\_urgent}; trực quan hóa trên Dashboard hai góc nhìn (Thực khách và Chủ quán); đồng thời cung cấp Trợ lý AI Chatbot 4 tầng Guardrails và công cụ tự động soạn thư phản hồi khách hàng đính kèm mã ưu đãi \texttt{CARE-*}. Hệ thống đã được kiểm thử tự động toàn diện qua 173 kịch bản kiểm thử đạt 100\% PASS.

% --- CÁC MỤC LỤC & DANH MỤC BẢNG BIỂU ---
\tableofcontents
\listoffigures
\listoftables

% =====================================================================
%  THÂN BÁO CÁO: 6 CHƯƠNG CHÍNH QUY (ĐÁNH SỐ Ả RẬP 1, 2, 3...)
% =====================================================================
\mainmatter
\dauPageStyle

% ==================== CHƯƠNG 1 ====================
\chapter{Tổng quan về Đề tài Nghiên cứu}

\section{Thông tin chung và Tính cấp thiết của Đề tài}
Ngành dịch vụ ẩm thực (Food and Beverage -- F\&B) tại Việt Nam có mức độ cạnh tranh vô cùng khốc liệt. Trước khi quyết định lựa chọn một địa điểm ăn uống, hơn 85\% thực khách hiện đại có thói quen tra cứu đánh giá trên Google Maps hoặc các chuyên trang ẩm thực như Foody.vn. Tuy nhiên, cả thực khách lẫn chủ nhà hàng đều đang đối mặt với bốn rào cản lớn:
\begin{enumerate}
    \item \textbf{Sự phân mảnh dữ liệu đa nền tảng (Data Silos):} Một nhà hàng lâu năm có thể sở hữu hàng trăm bình luận trên Foody.vn nhưng lại mới tạo địa điểm trên Google Maps, hoặc ngược lại, các quán ăn thế hệ mới có hàng ngàn đánh giá trên Google Maps nhưng rất thưa dữ liệu trên Foody. Nếu chỉ nhìn vào một nền tảng, kết quả phân tích sẽ bị lệch lạc nghiêm trọng.
    \item \textbf{Nhiễu loạn thông tin do đánh giá rác (Spam \& Noise Pollution):} Các trang đánh giá công cộng tại Việt Nam chứa lượng lớn bình luận rác như quảng cáo dịch vụ xe đưa đón, rao bán bất động sản, bình luận copy-paste để nhận xu thưởng (coin farming), hoặc các đoạn văn bản lỗi giao diện (\texttt{\{\{Model.Owner...\}\}}). Nếu đưa trực tiếp dữ liệu thô này vào mô hình AI, hệ thống vừa lãng phí chi phí API vừa làm sai lệch điểm số cảm xúc.
    \item \textbf{Điểm xếp hạng trung bình (Star Rating) không phản ánh đúng bản chất:} Một quán ăn đạt 4.0/5.0 sao không cho chủ quán biết khách hàng đang khen món ăn ngon hay đang phàn nàn gay gắt về thái độ nhân viên giữ xe và vấn đề vệ sinh.
    \item \textbf{Thiếu công cụ phản hồi và chăm sóc khách hàng tức thì:} Khi gặp một đánh giá 1 sao phàn nàn về trải nghiệm, chủ nhà hàng thường lúng túng hoặc phản hồi cảm tính, thiếu quy trình xin lỗi chuyên nghiệp kèm giải pháp đền bù (voucher) để giữ chân khách hàng.
\end{enumerate}

\section{Mục tiêu và Phạm vi của Hệ thống}
Hệ thống \textbf{Restaurant Intelligence} được thiết kế nhằm giải quyết trọn vẹn 4 bài toán trên thông qua quy trình xử lý thông minh khép kín:
\begin{itemize}
    \item \textbf{Thu thập đa nguồn \& Khớp nối thực thể:} Thu thập Foody.vn và Google Maps; so khớp mờ tiếng Việt ($\ge 85\%$), đối chiếu địa chỉ hành chính và mã băm SHA-256 khử trùng lặp.
    \item \textbf{Tiền xử lý lọc spam 11 trường hợp:} Lọc tự động các nhóm quảng cáo thuê xe, bất động sản, săn xu thưởng và dọn sạch mã lỗi AngularJS DOM.
    \item \textbf{Phân tích 5 khía cạnh ABSA qua Gemini API:} Bóc tách khía cạnh (\textit{Món ăn, Dịch vụ, Giá cả, Không gian, Vệ sinh/Vị trí}), nhãn cảm xúc, độ tự tin và cờ khẩn cấp \texttt{is\_urgent}.
    \item \textbf{Trợ lý AI Chatbot \& AI Review Responder:} Hỗ trợ thực khách hỏi đáp qua 4 tầng Guardrails và giúp chủ quán soạn thư xin lỗi/cảm ơn kèm cấp mã giảm giá \texttt{CARE-\{SHORTNAME\}-\{DISCOUNT\}}.
    \item \textbf{Dashboard Đa góc nhìn (Dual-Perspective BI):} Góc nhìn Thực khách (Món nên thử, so sánh đối đầu) và Chủ nhà hàng (Checklist vận hành, biểu đồ xu hướng, độ tin cậy mẫu $C(N)$).
\end{itemize}

\section{So sánh với các Giải pháp Hiện có}
Bảng~\ref{tab:comparison} đối chiếu sự khác biệt giữa việc tra cứu truyền thống trên các nền tảng hiện hành và hệ thống do nhóm phát triển:

\begin{table}[H]
\centering
\caption{So sánh tính năng giữa các nền tảng hiện có và Restaurant Intelligence}
\label{tab:comparison}
\vspace{0.15cm}
\footnotesize
\begin{tabularx}{\textwidth}{l c c c X}
\toprule
\rowcolor{blue!8}
\textbf{Tiêu chí kỹ thuật} & \textbf{Foody.vn} & \textbf{Google Maps} & \textbf{Social Listening} & \textbf{Restaurant Intelligence (Đề tài)} \\
\midrule
Nguồn dữ liệu & Đơn nguồn & Đơn nguồn & Mạng xã hội & \textbf{Hợp nhất Foody + Google Maps} \\
Lọc Spam đặc thù VN & Thấp & Trung bình & Theo từ khóa & \textbf{Bộ lọc đa tín hiệu 11 ca + dọn DOM lỗi} \\
Phân tích 5 khía cạnh & Điểm tĩnh & Từ khóa rời rạc & Tổng quan & \textbf{ABSA bằng Gemini + Cờ \texttt{is\_urgent}} \\
Đo độ tin cậy mẫu & Không có & Không có & Không có & \textbf{Hàm toán học logarit $C(N)$ minh bạch} \\
So sánh đối đầu 2 quán & Không có & Không có & Chi phí cao & \textbf{So sánh trực diện từng khía cạnh trên UI} \\
Trợ lý Q\&A \& Soạn Reply & Không có & Không có & Hạn chế & \textbf{Chatbot 4 tầng Guardrails + AI Responder} \\
\bottomrule
\end{tabularx}
\end{table}

% ==================== CHƯƠNG 2 ====================
\chapter{Cơ sở Lý thuyết và Mô hình Toán học}

\section{Bài toán Phân tích Cảm xúc theo Khía cạnh (ABSA)}
Khác với phân tích cảm xúc mức văn bản, bài toán \textbf{ABSA}~\cite{pontiki2014, zhang2022} giải quyết các câu nhận xét đa ý kiến. Bài toán được mô hình hóa thành phép ánh xạ từ văn bản $r_i$ sang tập hợp các bộ tứ (quadruples) theo Công thức~\eqref{eq:absa}:
\begin{equation}
\label{eq:absa}
    \mathcal{F}_{\text{ABSA}}(r_i) = \left\{ \left( a_k, s_k, c_k, u_k \right) \;\middle|\; a_k \in \mathcal{A},\; s_k \in \mathcal{S},\; c_k \in [0, 1],\; u_k \in \{0, 1\} \right\}
\end{equation}
Trong đó $\mathcal{A} = \{\texttt{food}, \texttt{service}, \texttt{price}, \texttt{ambiance}, \texttt{hygiene}\}$, $\mathcal{S} = \{\texttt{positive}, \texttt{neutral}, \texttt{negative}\}$, $c_k \in [0, 1]$ là độ tự tin và $u_k \in \{0, 1\}$ là cờ cảnh báo khẩn cấp (\texttt{is\_urgent}).

\section{Thuật toán Khớp nối Thực thể và Khử trùng lặp SHA-256}
Để nhận diện hai bản ghi từ Foody.vn và Google Maps thuộc cùng một nhà hàng mà không gộp nhầm các chi nhánh khác quận, nhóm áp dụng quy trình so khớp mờ theo Christen (2012)~\cite{christen2012}:
\begin{equation}
\label{eq:sim}
    \text{Sim}(\tilde{T}_F, \tilde{T}_G) = \left( 1 - \frac{d_{\text{Levenshtein}}(\tilde{T}_F, \tilde{T}_G)}{\max\left(|\tilde{T}_F|, |\tilde{T}_G|\right)} \right) \times 100\%
\end{equation}
Điều kiện ràng buộc gộp thực thể:
\begin{equation}
\label{eq:match}
    \text{Match}(R_F, R_G) = \mathbb{I}\Big( \text{Sim}(\tilde{T}_F, \tilde{T}_G) \ge 85\% \;\wedge\; \text{District}(A_F) \equiv \text{District}(A_G) \Big)
\end{equation}
Khử trùng lặp bình luận bằng chữ ký số SHA-256:
\begin{equation}
\label{eq:fingerprint}
    \text{Fingerprint}(r) = \text{SHA256}\Big( \texttt{restaurant\_id} \;\|\; \texttt{source} \;\|\; \text{NFKC}(\text{Lower}(\text{Trim}(r_{\text{text}}))) \Big)
\end{equation}

\section{Độ Tin cậy Cỡ mẫu Thống kê $C(N)$ và Điểm Sức khỏe}
Độ tin cậy cỡ mẫu tiệm cận logarit $C(N)$ dựa trên số lượng đánh giá hợp lệ $N$:
\begin{equation}
\label{eq:credibility}
    C(N) = \min\left( 1.0,\; \frac{\ln(N + 1)}{\ln(N_{\text{threshold}} + 1)} \right) \times 100\% \quad (N_{\text{threshold}} = 30)
\end{equation}
Điểm sức khỏe khía cạnh $k \in \mathcal{A}$ có trọng số theo độ tự tin $c_{i,k}$:
\begin{equation}
\label{eq:health_score}
    S_k = \frac{\sum_{i=1}^{M_k} w(s_{i,k}) \cdot c_{i,k}}{\sum_{i=1}^{M_k} c_{i,k}} \times 100 \quad \text{với } w(s) = \begin{cases} 1.0 & s = \texttt{positive} \\ 0.5 & s = \texttt{neutral} \\ 0.0 & s = \texttt{negative} \end{cases}
\end{equation}

% ==================== CHƯƠNG 3 ====================
\chapter{Phân tích và Thiết kế Hệ thống}

\section{Kiến trúc Hệ thống Tổng thể (3 Tầng)}
Hình~\ref{fig:arch_3layer} minh họa kiến trúc phân tầng của hệ thống gồm Tầng Trình diễn (React 19), Tầng Nghiệp vụ (FastAPI) và Tầng Dữ liệu \& AI (SQLite WAL + Gemini API).

\begin{figure}[H]
\centering
\resizebox{0.95\textwidth}{!}{%
\begin{tikzpicture}[
    layerbox/.style={draw=blue!75, fill=blue!3, dashed, thick, rounded corners=8pt, inner sep=14pt},
    comp/.style={draw=blue!85, fill=white, thick, rounded corners=5pt, minimum height=1.35cm, text width=5.6cm, align=center, font=\footnotesize},
    arrow/.style={-{Stealth[scale=0.95]}, thick, draw=gray!80!black}
]
    \node[comp] (react) at (-3.9, 0) {\textbf{React 19 Client Dashboard}\\(TypeScript, Recharts, Dual-View)};
    \node[comp] (aiwidget) at (3.9, 0) {\textbf{AI Assistant \& Modals}\\(Chatbot, Review Responder, URL Input)};
    
    \node[comp] (api) at (-3.9, -4.4) {\textbf{FastAPI Gateway \& Services}\\(REST API, BI Engine, Entity Match)};
    \node[comp] (crawler) at (3.9, -4.4) {\textbf{Multi-Source Crawler \& Filter}\\(Foody + Google Maps, Spam Filter)};

    \node[comp] (db) at (-3.9, -8.8) {\textbf{SQLite Database (WAL Mode)}\\(4 Bảng: Restaurants, Reviews, ABSA)};
    \node[comp] (gemini) at (3.9, -8.8) {\textbf{Google Gemini API Engine}\\(ABSA Batch, Q\&A, Auto-Reply)};

    \begin{scope}[on background layer]
        \node[layerbox, fit={(react) (aiwidget)}, label={[font=\bfseries\small, text=blue!80!black, yshift=2pt]above:Presentation Layer (Frontend SPA)}] (boxClient) {};
        \node[layerbox, fit={(api) (crawler)}, label={[font=\bfseries\small, text=blue!80!black, yshift=2pt]above:Business \& Processing Layer (FastAPI Backend)}] (boxBackend) {};
        \node[layerbox, fit={(db) (gemini)}, label={[font=\bfseries\small, text=blue!80!black, yshift=-2pt]below:Data \& Intelligence Layer}] (boxData) {};
    \end{scope}

    \draw[arrow] (react.south) -- node[left=4pt, pos=0.42, font=\scriptsize, align=right]{HTTP / JSON\\Request} (api.north);
    \draw[arrow] (aiwidget.south) -- ++(0, -1.05) -| ([xshift=1.4cm]api.north) node[pos=0.25, above, font=\scriptsize, fill=white, inner sep=1.5pt]{Ask / Reply / Analyze URL};
    \draw[arrow] (api.south) -- node[left=4pt, pos=0.45, font=\scriptsize]{SQLAlchemy ORM} (db.north);
    \draw[arrow] (crawler.south) -- node[right=4pt, pos=0.45, font=\scriptsize, align=left]{Batch Valid\\Reviews} (gemini.north);
    \draw[arrow] ([xshift=-1.4cm]crawler.south) -- ++(0, -0.95) -| ([xshift=1.2cm]db.north) node[pos=0.25, above, font=\scriptsize, fill=white, inner sep=1.5pt]{Lưu Reviews \& Cờ Spam / SHA-256};
    \draw[arrow] ([xshift=1.4cm]api.south) -- ++(0, -1.65) -| ([xshift=-1.4cm]gemini.north) node[pos=0.25, below, font=\scriptsize, fill=white, inner sep=1.5pt]{Rich Context Prompt};
    \draw[arrow] (gemini.west) -- node[above=2pt, font=\scriptsize]{Aspects / Sentiments} (db.east);
\end{tikzpicture}%
}
\caption{Kiến trúc tổng thể 3 tầng của hệ thống Restaurant Intelligence}
\label{fig:arch_3layer}
\end{figure}

\section{Sơ đồ Trường hợp Sử dụng (Use Case Diagram)}
Hình~\ref{fig:usecase_diagram} thể hiện sơ đồ Use Case chuẩn UML phân rã 4 tác nhân theo 2 cột chức năng đối xứng, triệt tiêu hoàn toàn lỗi chồng lấn đồ họa.

\begin{figure}[H]
\centering
\resizebox{0.96\textwidth}{!}{%
\begin{tikzpicture}[
    uc/.style={ellipse, draw=blue!80, fill=blue!6, thick, minimum height=1.0cm, text width=3.8cm, align=center, font=\scriptsize, inner sep=2pt},
    uc_sub/.style={ellipse, draw=teal!80, fill=teal!6, thick, minimum height=0.9cm, text width=3.4cm, align=center, font=\tiny, inner sep=1.5pt},
    sysbox/.style={draw=gray!70, dashed, thick, inner sep=14pt, rounded corners=8pt},
    actor/.style={font=\bfseries\scriptsize, align=center},
    rel/.style={-{Stealth[scale=0.85]}, thick},
    dep/.style={-{Stealth[scale=0.8]}, dashed, thick, draw=blue!75!black, font=\tiny}
]
    % --- ACTORS TRÁI ---
    \node[actor] (dinerLabel) at (-7.0, 3.2) {Thực khách\\(Diner)};
    \draw[thick] (-7.0, 2.5) circle (0.3cm);
    \draw[thick] (-7.0, 2.2) -- (-7.0, 1.4);
    \draw[thick] (-7.5, 1.85) -- (-6.5, 1.85);
    \draw[thick] (-7.0, 1.4) -- (-7.4, 0.7);
    \draw[thick] (-7.0, 1.4) -- (-6.6, 0.7);
    \coordinate (diner) at (-6.3, 1.85);

    \node[actor] (ownerLabel) at (-7.0, -1.8) {Chủ nhà hàng\\(Owner)};
    \draw[thick] (-7.0, -2.5) circle (0.3cm);
    \draw[thick] (-7.0, -2.8) -- (-7.0, -3.6);
    \draw[thick] (-7.5, -3.15) -- (-6.5, -3.15);
    \draw[thick] (-7.0, -3.6) -- (-7.4, -4.3);
    \draw[thick] (-7.0, -3.6) -- (-6.6, -4.3);
    \coordinate (owner) at (-6.3, -3.15);

    % --- CỘT 1: THỰC KHÁCH (x = -2.1) ---
    \node[uc] (uc1) at (-2.1, 4.0) {1. Tìm kiếm \& Lọc quán\\(63 tỉnh thành)};
    \node[uc] (uc2) at (-2.1, 2.2) {2. Xem Dashboard phân tích\\(Góc nhìn Thực khách)};
    \node[uc] (uc3) at (-2.1, 0.4) {3. So sánh đối đầu\\(2 nhà hàng trực diện)};
    \node[uc] (uc4) at (-2.1, -1.4) {4. Lọc đánh giá theo\\5 khía cạnh \& nhãn Spam};
    \node[uc] (uc5) at (-2.1, -3.2) {5. Hỏi đáp Trợ lý AI Chatbot\\(4 tầng Guardrails)};

    % --- CỘT 2: VẬN HÀNH & HỆ THỐNG (x = 2.4) ---
    \node[uc] (uc6) at (2.4, 4.0) {6. Xem Dashboard Chuyên sâu\\(SWOT, Độ tin cậy $C(N)$)};
    \node[uc] (uc7) at (2.4, 2.2) {7. Xem Checklist Vận hành\\\& Cảnh báo \texttt{is\_urgent}};
    \node[uc] (uc8) at (2.4, 0.4) {8. Soạn thư phản hồi AI\\(4 tông giọng)};
    \node[uc_sub] (uc8_ext) at (2.4, -1.2) {8.1. Đính kèm mã Voucher\\(\texttt{CARE-*})};
    \node[uc] (uc9) at (2.4, -2.8) {9. Yêu cầu phân tích URL mới\\(Live Foody / Maps)};
    \node[uc_sub] (uc_crawl) at (2.4, -4.4) {10. Pipeline cào đa nguồn,\\Lọc Spam 11 ca \& ABSA};

    % --- ACTORS PHẢI ---
    \node[draw=purple!80, fill=purple!8, thick, rounded corners=4pt, font=\scriptsize\bfseries, align=center, text width=2.4cm] (gemini) at (6.8, -0.4) {Google Gemini\\API Engine\\(LLM Cloud)};

    \node[actor] (cronLabel) at (6.8, -3.6) {Hệ thống tự động\\(Auto-Enrich Cron)};
    \draw[thick] (6.8, -4.2) circle (0.3cm);
    \draw[thick] (6.8, -4.5) -- (6.8, -5.3);
    \draw[thick] (6.3, -4.85) -- (7.3, -4.85);
    \draw[thick] (6.8, -5.3) -- (6.4, -6.0);
    \draw[thick] (6.8, -5.3) -- (7.2, -6.0);
    \coordinate (cron) at (6.3, -4.85);

    % --- SYSTEM BOUNDARY ---
    \node[sysbox, fit={(uc1) (uc5) (uc6) (uc_crawl)}, label={[font=\bfseries\small, text=blue!85!black, yshift=4pt]above:Hệ thống Restaurant Intelligence}] (sysContainer) {};

    % --- LIÊN KẾT TÁC NHÂN ---
    \draw[rel] (diner) -- (uc1.west);
    \draw[rel] (diner) -- (uc2.west);
    \draw[rel] (diner) -- (uc3.west);
    \draw[rel] (diner) -- (uc4.west);
    \draw[rel] (diner) -- (uc5.west);

    \draw[rel] (owner) -- (uc2.south west);
    \draw[rel] (owner) -- (uc6.west);
    \draw[rel] (owner) -- (uc7.west);
    \draw[rel] (owner) -- (uc8.west);
    \draw[rel] (owner) -- (uc9.west);

    \draw[rel] (cron) -- (uc_crawl.east);

    % --- INCLUDE & EXTEND ---
    \draw[dep] (uc8) -- node[right=1pt, font=\tiny]{\textit{<<extend>>}} (uc8_ext);
    \draw[dep] (uc9) -- node[right=1pt, font=\tiny]{\textit{<<include>>}} (uc_crawl);

    \draw[dep] (uc5.east) -- node[above, font=\tiny, sloped]{\textit{<<include>>}} (gemini.west);
    \draw[dep] (uc8.east) -- node[above, font=\tiny, sloped]{\textit{<<include>>}} (gemini.west);
    \draw[dep] (uc_crawl.north east) -- node[below, font=\tiny, sloped]{\textit{<<include>>}} (gemini.south);
\end{tikzpicture}%
}
\caption{Sơ đồ Use Case tổng thể được phân rã theo 4 tác nhân và 2 cột chức năng đối xứng}
\label{fig:usecase_diagram}
\end{figure}

\begin{table}[H]
\centering
\caption{Đặc tả phân quyền và vai trò của các Use Case trọng tâm}
\label{tab:usecase_spec}
\vspace{0.15cm}
\footnotesize
\begin{tabularx}{\textwidth}{l l l X}
\toprule
\rowcolor{blue!8}
\textbf{Mã UC} & \textbf{Tên Use Case} & \textbf{Tác nhân chính} & \textbf{Mô tả luồng nghiệp vụ và giá trị đầu ra} \\
\midrule
\textbf{UC-01} & Tìm kiếm \& Lọc quán & Thực khách / Chủ quán & Lọc 50 nhà hàng theo 63 tỉnh/thành, loại hình ẩm thực và khoảng giá. \\
\textbf{UC-02} & Dashboard Thực khách & Thực khách & Xem tổng quan điểm 5 khía cạnh, món ăn được khen/chê và trích dẫn nổi bật. \\
\textbf{UC-03} & So sánh đối đầu 2 quán & Thực khách & Đặt 2 quán cạnh nhau so sánh chi tiết điểm số, tỷ lệ spam và thế mạnh. \\
\textbf{UC-05} & Hỏi đáp Trợ lý AI & Thực khách & Chatbot giải đáp thắc mắc dựa trên dữ liệu thật qua 4 tầng rào chắn Guardrails. \\
\textbf{UC-06} & Dashboard Chuyên sâu & Chủ nhà hàng & Đánh giá xu hướng thời gian, phân tích SWOT và độ tin cậy mẫu thống kê $C(N)$. \\
\textbf{UC-07} & Checklist \& Cảnh báo & Chủ nhà hàng & Theo dõi các đầu việc vận hành cần sửa chữa và phát hiện các bài đánh giá \texttt{is\_urgent}. \\
\textbf{UC-08} & Soạn thư phản hồi AI & Chủ nhà hàng & Chọn tông giọng; AI soạn thư xin lỗi/cảm ơn kèm mở rộng (\textit{<<extend>>}) cấp mã \texttt{CARE-*}. \\
\textbf{UC-09} & Phân tích URL Mới & Chủ nhà hàng & Dán link Foody/Maps; tự động gọi (\textit{<<include>>}) pipeline cào, lọc spam và phân tích ngay. \\
\bottomrule
\end{tabularx}
\end{table}

\section{Sơ đồ Luồng Hoạt động (Activity Diagram)}
Hình~\ref{fig:activity_diagram} mô tả quy trình thu thập, lọc rác, gom lô batch 20 và phân tích ABSA.

\begin{figure}[H]
\centering
\resizebox{0.78\textwidth}{!}{%
\begin{tikzpicture}[
    node distance=0.65cm,
    startstop/.style={draw, circle, fill=black, inner sep=4pt},
    endnode/.style={draw, double, circle, fill=black, inner sep=3.2pt},
    action/.style={draw=blue!80, fill=blue!8, thick, rounded corners, minimum width=5.6cm, minimum height=0.75cm, align=center, font=\footnotesize},
    decision/.style={draw=orange!90, fill=orange!10, diamond, aspect=2.2, thick, font=\scriptsize, align=center, inner sep=2pt},
    arrow/.style={-{Stealth[scale=0.95]}, thick}
]
    \node[startstop] (start) {};
    \node[action, below=of start] (crawl) {Crawl đánh giá từ Foody.vn / Google Maps};
    \node[action, below=of crawl] (dedup) {Khớp thực thể quán \& Khử trùng lặp SHA-256};
    \node[action, below=of dedup] (filterSpam) {Kiểm tra \& Đánh dấu Spam 11 ca (\texttt{ReviewFilter})};
    \node[decision, below=of filterSpam] (isSpam) {Đánh giá là Spam\\hoặc Quá ngắn?};
    \node[action, right=1.2cm of isSpam, minimum width=4.4cm, text width=4.2cm] (markOnly) {Gán cờ \texttt{is\_spam=1, is\_analyzed=1}\\(Bỏ qua gọi Gemini AI)};
    \node[action, below=0.85cm of isSpam] (cleanValid) {Đưa vào hàng đợi phân tích hợp lệ};
    \node[decision, below=of cleanValid] (batchCheck) {Đủ kích thước lô\\(\texttt{BATCH\_SIZE = 20})?};
    \node[action, below=0.85cm of batchCheck] (callAI) {Gọi Gemini API phân tích 5 khía cạnh \& \texttt{is\_urgent}};
    \node[action, below=of callAI] (saveDB) {Lưu vào SQLite \& Tính điểm tin cậy mẫu $C(N)$};
    \node[action, below=of saveDB] (visualize) {Xóa Cache RAM, cập nhật Dashboard \& Trợ lý AI};
    \node[endnode, below=of visualize] (end) {};

    \draw[arrow] (start) -- (crawl);
    \draw[arrow] (crawl) -- (dedup);
    \draw[arrow] (dedup) -- (filterSpam);
    \draw[arrow] (filterSpam) -- (isSpam);
    \draw[arrow] (isSpam) -- node[above, font=\scriptsize]{Có} (markOnly);
    \draw[arrow] (isSpam) -- node[right=2pt, font=\scriptsize]{Không} (cleanValid);
    \draw[arrow] (markOnly.south) |- (saveDB.east);
    \draw[arrow] (cleanValid) -- (batchCheck);
    \draw[arrow] (batchCheck) -- node[right=2pt, font=\scriptsize]{Đủ batch} (callAI);
    \draw[arrow] (batchCheck.west) -- ++(-1.5,0) |- node[pos=0.25, left, font=\scriptsize]{Gom tiếp} (cleanValid.west);
    \draw[arrow] (callAI) -- (saveDB);
    \draw[arrow] (saveDB) -- (visualize);
    \draw[arrow] (visualize) -- (end);
\end{tikzpicture}%
}
\caption{Sơ đồ Activity thể hiện luồng xử lý dữ liệu và tối ưu hóa chi phí gọi AI}
\label{fig:activity_diagram}
\end{figure}

\section{Sơ đồ Tuần tự (Sequence Diagram)}
Hình~\ref{fig:sequence_diagram} minh họa thông điệp trao đổi giữa User/React UI, FastAPI, Gemini API và SQLite DB.

\begin{figure}[H]
\centering
\resizebox{0.95\textwidth}{!}{%
\begin{tikzpicture}[
    lifeline/.style={draw=gray!60, dashed, thick},
    obj/.style={draw=blue!70, fill=blue!10, thick, rounded corners, minimum width=2.5cm, minimum height=0.85cm, font=\footnotesize\bfseries, align=center},
    msg/.style={-{Stealth[scale=0.9]}, thick, font=\scriptsize},
    returnmsg/.style={-{Stealth[scale=0.9]}, dashed, thick, font=\scriptsize}
]
    \node[obj] (client) at (0,0) {User / React UI};
    \node[obj] (api) at (4.3,0) {FastAPI Backend};
    \node[obj] (ai) at (8.6,0) {Gemini API};
    \node[obj] (db) at (12.9,0) {SQLite DB};

    \draw[lifeline] (client) -- ++(0, -10.6);
    \draw[lifeline] (api) -- ++(0, -10.6);
    \draw[lifeline] (ai) -- ++(0, -10.6);
    \draw[lifeline] (db) -- ++(0, -10.6);

    \draw[msg] (0, -1.0) -- node[above, fill=white, inner sep=1pt]{1. Yêu cầu Dashboard nhà hàng} (4.3, -1.0);
    \draw[msg] (4.3, -1.8) -- node[above, fill=white, inner sep=1pt]{2. Truy vấn thống kê, khía cạnh \& đánh giá} (12.9, -1.8);
    \draw[returnmsg] (12.9, -2.6) -- node[above, fill=white, inner sep=1pt]{3. Trả về reviews (kèm nhãn Spam \& ABSA)} (4.3, -2.6);
    \draw[returnmsg] (4.3, -3.4) -- node[above, fill=white, inner sep=1pt]{4. Tổng hợp SWOT, Checklist \& hiển thị UI} (0, -3.4);

    \draw[msg] (0, -4.5) -- node[above, fill=white, inner sep=1pt]{5. Gửi câu hỏi Chatbot / Yêu cầu soạn Reply} (4.3, -4.5);
    \draw[msg] (4.3, -5.3) -- node[above, fill=white, inner sep=1pt]{6. Truy xuất ngữ cảnh 5 khía cạnh \& 12 trích dẫn} (12.9, -5.3);
    \draw[returnmsg] (12.9, -6.1) -- node[above, fill=white, inner sep=1pt]{7. Dữ liệu ngữ cảnh nhà hàng thực tế} (4.3, -6.1);
    \draw[msg] (4.3, -7.0) -- node[above, fill=white, inner sep=1pt]{8. Gửi System Instruction + XML Prompt} (8.6, -7.0);
    \draw[returnmsg] (8.6, -7.9) -- node[above, fill=white, inner sep=1pt]{9. Trả lời phân tích / Thư phản hồi + Mã Voucher} (4.3, -7.9);
    \draw[msg] (4.3, -8.6) -- ++(1.2, 0) -- ++(0, -0.55) -- node[right=1.2cm, pos=0.5, font=\scriptsize, fill=white, inner sep=1pt]{10. Lưu In-Memory Cache (0ms lần sau)} (4.3, -9.15);
    \draw[returnmsg] (4.3, -9.9) -- node[above, fill=white, inner sep=1pt]{11. Hiển thị câu trả lời AI / Bản thảo thư} (0, -9.9);
\end{tikzpicture}%
}
\caption{Sơ đồ tuần tự luồng tải Dashboard, hỏi đáp AI Chatbot và soạn thư phản hồi}
\label{fig:sequence_diagram}
\end{figure}

\section{Sơ đồ Thực thể Liên kết (ERD) và Luồng Dữ liệu (DFD)}
Cấu trúc cơ sở dữ liệu quan hệ được mô tả tại Hình~\ref{fig:erd_diagram}, trong khi luồng thông tin DFD Mức 0 và Mức 1 được minh họa tại Hình~\ref{fig:dfd_level0} và Hình~\ref{fig:dfd_level1}.

\begin{figure}[H]
\centering
\resizebox{0.95\textwidth}{!}{%
\begin{tikzpicture}[
    entity/.style={draw=blue!80, fill=blue!10, rectangle, thick, minimum width=2.9cm, minimum height=0.95cm, align=center, font=\small\bfseries},
    attr/.style={draw=black!70, fill=white, ellipse, font=\scriptsize, align=center, inner sep=2pt},
    rel/.style={draw=orange!80, fill=orange!10, diamond, aspect=1.8, thick, font=\scriptsize\bfseries, align=center}
]
    \node[entity] (rest) at (0, 0) {Restaurant};
    \node[rel] (has) at (4.4, 0) {Has};
    \node[entity] (rev) at (8.8, 0) {Review};
    \node[rel] (produces) at (8.8, -3.2) {Produces};
    \node[entity] (aspect) at (8.8, -6.4) {ReviewAnalysis};
    \node[entity] (queue) at (0, -5.4) {CrawlRequest};

    \draw[thick] (rest) -- node[above, font=\small]{1} (has);
    \draw[thick] (has) -- node[above, font=\small]{N} (rev);
    \draw[thick] (rev) -- node[right, font=\small]{1} (produces);
    \draw[thick] (produces) -- node[right, font=\small]{N} (aspect);

    \node[attr] (r_id) at (-2.7, 1.5) {\underline{id}};
    \node[attr] (r_name) at (-0.6, 2.0) {name};
    \node[attr] (r_city) at (1.8, 1.6) {city / address};
    \node[attr] (r_urls) at (-3.1, -0.1) {foody / gmaps\_url};
    \node[attr] (r_cred) at (-1.8, -1.7) {credibility\_score};
    \draw (rest) -- (r_id); \draw (rest) -- (r_name); \draw (rest) -- (r_city);
    \draw (rest) -- (r_urls); \draw (rest) -- (r_cred);

    \node[attr] (rv_id) at (6.5, 1.6) {\underline{id}};
    \node[attr] (rv_src) at (8.6, 2.0) {source};
    \node[attr] (rv_text) at (10.9, 1.6) {text / rating};
    \node[attr] (rv_fp) at (12.1, 0.2) {fingerprint};
    \node[attr] (rv_spam) at (11.9, -1.3) {\textbf{is\_spam / category}};
    \draw (rev) -- (rv_id); \draw (rev) -- (rv_src); \draw (rev) -- (rv_text);
    \draw (rev) -- (rv_fp); \draw (rev) -- (rv_spam);

    \node[attr] (a_id) at (5.6, -5.6) {\underline{id}};
    \node[attr] (a_cat) at (5.6, -7.0) {aspect};
    \node[attr] (a_pol) at (7.8, -8.2) {sentiment};
    \node[attr] (a_conf) at (10.3, -8.2) {confidence};
    \node[attr] (a_urg) at (12.0, -6.4) {\textbf{is\_urgent}};
    \draw (aspect) -- (a_id); \draw (aspect) -- (a_cat); \draw (aspect) -- (a_pol);
    \draw (aspect) -- (a_conf); \draw (aspect) -- (a_urg);

    \node[attr] (q_id) at (-2.7, -4.4) {\underline{id}};
    \node[attr] (q_query) at (-2.4, -6.4) {query / city};
    \node[attr] (q_stat) at (0.4, -7.2) {status};
    \draw (queue) -- (q_id); \draw (queue) -- (q_query); \draw (queue) -- (q_stat);
\end{tikzpicture}%
}
\caption{Sơ đồ Thực thể Liên kết (ERD) của cơ sở dữ liệu Restaurant Intelligence}
\label{fig:erd_diagram}
\end{figure}

\begin{figure}[H]
\centering
\resizebox{0.86\textwidth}{!}{%
\begin{tikzpicture}[
    ext/.style={draw=black!80, fill=gray!10, rectangle, thick, minimum width=3.1cm, minimum height=1.2cm, align=center, font=\footnotesize\bfseries},
    proc/.style={draw=blue!80, fill=blue!10, circle, thick, minimum size=3.0cm, align=center, font=\footnotesize\bfseries},
    arrow/.style={-{Stealth[scale=0.95]}, thick, font=\scriptsize}
]
    \node[proc] (p0) at (0,0) {0.0\\Hệ thống\\Restaurant\\Intelligence};
    \node[ext] (foody) at (-6.2,0) {Foody.vn \&\\Google Maps};
    \node[ext] (user) at (6.2,0) {Thực khách /\\Chủ nhà hàng};

    \draw[arrow] (foody.east) -- node[above=3pt, align=center]{Đánh giá \&\\Thông tin quán} (p0.west);
    \draw[arrow] (user.north west) to[bend right=24] node[above=3pt]{Tìm kiếm, Dán URL \& Hỏi AI} (p0.north east);
    \draw[arrow] (p0.south east) to[bend right=24] node[below=3pt]{Dashboard, Cảnh báo \& Thư Reply} (user.south west);
\end{tikzpicture}%
}
\caption{Sơ đồ Luồng dữ liệu Mức 0 (DFD Context Diagram)}
\label{fig:dfd_level0}
\end{figure}

\begin{figure}[H]
\centering
\resizebox{0.88\textwidth}{!}{%
\begin{tikzpicture}[
    proc/.style={draw=blue!80, fill=blue!10, rounded corners=10pt, thick, text width=3.5cm, minimum height=1.15cm, align=center, font=\scriptsize\bfseries},
    store/.style={draw=black!80, fill=yellow!10, thick, text width=3.8cm, minimum height=0.9cm, align=center, font=\scriptsize},
    arrow/.style={-{Stealth[scale=0.95]}, thick, font=\scriptsize}
]
    \node[proc] (p1) at (-3.8, 3.6) {1.0 Thu thập đa nguồn\\\& Lọc Spam 11 ca};
    \node[proc] (p2) at (3.8, 3.6) {2.0 Phân tích ABSA\\(Google Gemini AI)};

    \node[store] (d1) at (-3.8, 0.8) {[D1] Bảng \texttt{restaurants}\\và \texttt{reviews}};
    \node[proc] (p3) at (3.8, 0.8) {3.0 Tổng hợp BI\\\& Dashboard Kép};

    \node[proc] (p4) at (-3.8, -2.0) {4.0 Trợ lý Chatbot \&\\AI Review Responder};
    \node[store] (d2) at (3.8, -2.0) {[D2] Bảng kết quả\\\texttt{review\_analysis}};

    \draw[arrow] (p1.south) -- node[left=4pt]{Lưu review \& cờ Spam} (d1.north);
    \draw[arrow] (d1.east) -- ++(1.3, 0) |- node[pos=0.72, above=2pt]{Review hợp lệ (Batch 20)} (p2.west);
    \draw[arrow] (p2.south) -- node[right=4pt]{Cập nhật tức thì} (p3.north);
    \draw[arrow] (p2.east) -- ++(1.1, 0) |- node[pos=0.25, right=3pt, align=left]{Lưu kết quả\\5 khía cạnh} (d2.east);
    \draw[arrow] (d2.north) -- node[left=4pt]{Đọc 5 khía cạnh} (p3.south);
    \draw[arrow] (d1.south) -- node[left=4pt]{Trích dẫn mẫu} (p4.north);
    \draw[arrow] (d2.west) -- node[below=3pt]{Dữ liệu thống kê ABSA} (p4.east);
    \draw[arrow] ([yshift=8pt]p4.east) -- ++(1.6, 0) |- node[pos=0.22, left=2pt, align=right]{Tư vấn \& Thư trả lời\\kèm Mã Voucher} ([yshift=-8pt]p3.west);
\end{tikzpicture}%
}
\caption{Sơ đồ Luồng dữ liệu Mức 1 (DFD Level 1)}
\label{fig:dfd_level1}
\end{figure}

% ==================== CHƯƠNG 4 ====================
\chapter{Hiện thực và Cài đặt Hệ thống}

\section{Hệ sinh thái Công nghệ Sử dụng}
Chi tiết các thư viện và công nghệ cốt lõi được liệt kê trong Bảng~\ref{tab:tech_stack}:

\begin{table}[H]
\centering
\caption{Đặc tả chi tiết ngăn xếp công nghệ (Technology Stack) của dự án}
\label{tab:tech_stack}
\vspace{0.15cm}
\footnotesize
\begin{tabularx}{\textwidth}{l l X}
\toprule
\rowcolor{blue!8}
\textbf{Phân hệ} & \textbf{Công nghệ / Thư viện} & \textbf{Vai trò và Lý do lựa chọn trong hệ thống} \\
\midrule
\textbf{Multi-Source Crawler} & Playwright (Chromium) & Thu thập Google Maps đa chiều, vượt lazy-loading DOM bằng container ảo. \\
 & Selenium \& BS4 & Thu thập Foody.vn kết hợp Fast HTTP Inspector ($<0.5\,\mathrm{s}$). \\
\midrule
\textbf{NLP \& Tiền xử lý} & Regex \& Unicodedata & Chuẩn hóa Unicode NFKC/NFD, vận hành bộ lọc Spam 11 trường hợp và dọn DOM. \\
 & Difflib / Entity Matcher & So khớp mờ tên quán và đối chiếu địa chỉ cấp Quận/Huyện ($\ge 85\%$). \\
\midrule
\textbf{Trí tuệ nhân tạo (LLM)} & Google GenAI SDK~\cite{google2025} & Kết nối \texttt{gemini-3.5-flash-lite} \& \texttt{gemini-2.5-flash} có Structured JSON Schema. \\
\midrule
\textbf{Cơ sở dữ liệu} & SQLite 3 (WAL Mode)~\cite{hipp2024} & Ghi nhật ký trước giúp đọc/ghi đồng thời không khóa bảng (\texttt{database is locked}). \\
 & SQLAlchemy 2.0 ORM & Quản lý 4 bảng quan hệ, bảo vệ chống SQL Injection tuyệt đối. \\
\midrule
\textbf{Backend \& Bảo mật} & FastAPI \& Uvicorn~\cite{ramirez2025} & RESTful API bất đồng bộ (ASGI), Pydantic v2 validation, 10 lớp bảo mật. \\
\midrule
\textbf{Frontend SPA} & React 19 \& TypeScript & Dashboard tốc độ cao với Vite, kiểm tra kiểu tĩnh an toàn. \\
 & Tailwind v4 \& Recharts & Trực quan hóa biểu đồ cột chồng 5 khía cạnh, xu hướng thời gian, thẻ so sánh. \\
\midrule
\textbf{Kiểm thử \& DevOps} & Pytest \& Docker & 173 kịch bản kiểm thử tự động, Docker Multi-stage, GitHub Actions CI/CD. \\
\bottomrule
\end{tabularx}
\end{table}

\section{Đặc tả Cơ sở Dữ liệu và RESTful API Endpoints}
Cơ sở dữ liệu lưu trữ trên SQLite WAL mode gồm bảng \texttt{restaurants}, \texttt{reviews} và \texttt{review\_analysis}. Bảng~\ref{tab:api_endpoints} liệt kê các endpoint chính:

\begin{table}[H]
\centering
\caption{Danh mục các RESTful API Endpoints cốt lõi trên FastAPI Backend}
\label{tab:api_endpoints}
\vspace{0.15cm}
\footnotesize
\begin{tabularx}{\textwidth}{l l X}
\toprule
\rowcolor{blue!8}
\textbf{Phương thức} & \textbf{Đường dẫn Endpoint} & \textbf{Chức năng xử lý và Cơ chế bảo vệ} \\
\midrule
\texttt{GET} & \texttt{/api/restaurants} & Danh sách quán có phân trang, tìm kiếm tên, lọc tỉnh thành và sắp xếp. \\
\texttt{GET} & \texttt{/api/restaurants/\{id\}/dashboard} & Toàn bộ gói Dashboard Kép: 5 khía cạnh, xu hướng, SWOT, Checklist vận hành. \\
\texttt{GET} & \texttt{/api/restaurants/compare} & Đối đầu trực tiếp 2 quán trên cả 5 khía cạnh và tỷ lệ spam. \\
\texttt{POST} & \texttt{/api/crawl/analyze-url} & Chống SSRF, cào link thời gian thực (\texttt{Semaphore(2)}), lọc spam và phân tích AI. \\
\texttt{POST} & \texttt{/api/assistant/ask} & Trợ lý AI Chatbot: Kiểm tra 4 tầng Guardrails, nạp 12 trích dẫn thật vào XML Prompt. \\
\texttt{POST} & \texttt{/api/reviews/\{id\}/reply} & AI Review Responder: Soạn thư theo 4 tông giọng kèm cấp mã voucher \texttt{CARE-*}. \\
\bottomrule
\end{tabularx}
\end{table}

\section{Bộ lọc Spam Đa tín hiệu 11 Trường hợp (\texttt{ReviewFilter})}
Bảng~\ref{tab:spam_cases} mô tả 11 kịch bản phát hiện spam và bảo vệ ngữ cảnh ẩm thực:

\begin{table}[H]
\centering
\caption{Đặc tả 11 kịch bản nhận diện Spam và Bảo vệ ngữ cảnh trong \texttt{ReviewFilter}}
\label{tab:spam_cases}
\vspace{0.15cm}
\scriptsize
\begin{tabularx}{\textwidth}{c l l X}
\toprule
\rowcolor{blue!8}
\textbf{STT} & \textbf{Mã kịch bản} & \textbf{Hành động} & \textbf{Quy tắc nhận diện (Regex / Heuristics) \& Ví dụ thực tế} \\
\midrule
1 & \texttt{valid\_normal} & Giữ lại (\texttt{is\_spam=0}) & Đánh giá thực khách bình thường có nhận xét cụ thể về món ăn, dịch vụ, không gian. \\
2 & \texttt{low\_information} & Đánh dấu (\texttt{is\_spam=1}) & Bình luận quá ngắn ($< 10$ ký tự hoặc $< 3$ từ), chỉ chứa biểu tượng cảm xúc (\textit{``ok'', ``...''}). \\
3 & \texttt{car\_rental\_spam} & Đánh dấu (\texttt{is\_spam=1}) & Quảng cáo dịch vụ xe đưa đón, taxi, xe ghép tuyến (\textit{``xe 4-7-16 chỗ đi Bà Nà, Hội An...''}). \\
4 & \texttt{contact\_ad\_spam} & Đánh dấu (\texttt{is\_spam=1}) & Chứa số điện thoại rác đi kèm từ khóa mời chào (\textit{``liên hệ Zalo'', ``inbox hotline''}). \\
5 & \texttt{context\_gia\_re} & \textbf{Bảo vệ (\texttt{is\_spam=0})} & \textbf{Luật chống lọc nhầm:} Cụm từ \textit{``giá rẻ''} đi kèm món ăn (\textit{``hải sản tươi giá rẻ''}). \\
6 & \texttt{context\_booking\_phone} & \textbf{Bảo vệ (\texttt{is\_spam=0})} & \textbf{Luật chống lọc nhầm:} Bình luận chia sẻ SĐT đặt bàn của quán (\textit{``gọi số quán để đặt trước''}). \\
7 & \texttt{duplicate\_spam} & Đánh dấu (\texttt{is\_spam=1}) & Bình luận trùng khớp mã băm SHA-256 trên nhiều quán khác nhau. \\
8 & \texttt{real\_estate\_finance} & Đánh dấu (\texttt{is\_spam=1}) & Rao vặt bán đất nền, căn hộ, vay vốn tín chấp, tuyển cộng tác viên tại nhà. \\
9 & \texttt{name\_echo\_bot} & Đánh dấu (\texttt{is\_spam=1}) & Bot tự động chỉ copy nguyên văn tên nhà hàng và địa chỉ vào phần bình luận. \\
10 & \texttt{generic\_template\_bot} & Đánh dấu (\texttt{is\_spam=1}) & Câu văn sáo rỗng hoàn toàn không có danh từ chỉ món ăn hay dịch vụ. \\
11 & \texttt{coin\_farming} & Đánh dấu (\texttt{is\_spam=1}) & Bình luận cày xu trên app (\textit{``viết cho đủ 50 ký tự để nhận xu thưởng''}) kèm ký tự lặp. \\
\bottomrule
\end{tabularx}
\end{table}

\section{Kiến trúc Gemini AI, 4 Tầng Guardrails và AI Review Responder}
Hình~\ref{fig:ai_pipeline} minh họa kiến trúc 4 tầng bảo vệ Guardrails ngăn chặn hoàn toàn tấn công Prompt Injection và hiện tượng ảo giác (Zero Hallucination).

\begin{figure}[H]
\centering
\resizebox{0.95\textwidth}{!}{%
\begin{tikzpicture}[
    guard/.style={draw=teal!80!black, fill=teal!6, thick, rounded corners=4pt, minimum width=11.8cm, minimum height=0.9cm, align=center, font=\footnotesize},
    model/.style={draw=blue!80, fill=blue!8, thick, rounded corners=5pt, minimum width=3.6cm, minimum height=1.1cm, align=center, font=\scriptsize\bfseries},
    cache/.style={draw=purple!80, fill=purple!8, thick, rounded corners=5pt, minimum width=4.8cm, minimum height=0.95cm, align=center, font=\footnotesize\bfseries},
    arrow/.style={-{Stealth[scale=0.95]}, thick}
]
    \node[guard] (g1) at (0, 0) {\textbf{Tầng 1 -- Input Sanitization \& Length Guard:} Kiểm tra độ dài ($\le 500$ ký tự), lọc thẻ điều khiển độc hại};
    \node[guard] (g2) at (0, -1.35) {\textbf{Tầng 2 -- Domain Intent Classifier:} Từ chối câu hỏi ngoài phạm vi ẩm thực / nhà hàng (code, chính trị, toán...)};
    \node[guard] (g3) at (0, -2.7) {\textbf{Tầng 3 -- XML Delimiter Sandboxing:} Bọc câu hỏi trong thẻ \texttt{<user\_question>...</user\_question>} tách biệt lệnh hệ thống};
    \node[guard] (g4) at (0, -4.05) {\textbf{Tầng 4 -- Grounded Context Injection:} Chỉ nạp thống kê 5 khía cạnh thật + 12 trích dẫn từ SQLite DB};

    \node[cache] (ttl) at (0, -5.55) {Kiểm tra In-Memory TTL Cache (Trả kết quả $0\,\mathrm{ms}$ nếu trùng khóa)};

    \node[model] (m1) at (-4.2, -7.3) {Mô hình ưu tiên 1:\\\texttt{gemini-3.5-flash-lite}\\(Tốc độ cao, Quota lớn)};
    \node[model] (m2) at (0, -7.3) {Mô hình dự phòng 2:\\\texttt{gemini-2.5-flash}\\(Khi gặp lỗi 429 / 503)};
    \node[model, draw=orange!85!black, fill=orange!10] (m3) at (4.2, -7.3) {Dự phòng cuối (Local):\\\textbf{Deterministic Template}\\(100\% Uptime khi mất mạng)};

    \draw[arrow] (g1) -- (g2);
    \draw[arrow] (g2) -- (g3);
    \draw[arrow] (g3) -- (g4);
    \draw[arrow] (g4) -- (ttl);
    \draw[arrow] (ttl.south) -- ++(0, -0.35) -| (m1.north);
    \draw[arrow, dashed, draw=red!80!black] (m1.east) -- node[above, font=\tiny]{Lỗi 429/Quota} (m2.west);
    \draw[arrow, dashed, draw=red!80!black] (m2.east) -- node[above, font=\tiny]{Timeout/Offline} (m3.west);
\end{tikzpicture}%
}
\caption{Kiến trúc 4 tầng bảo vệ Guardrails và cơ chế chuyển đổi dự phòng Cascading Fallback}
\label{fig:ai_pipeline}
\end{figure}

Chủ nhà hàng có thể chọn 4 tông giọng: \textit{Chuyên nghiệp (\texttt{professional})}, \textit{Thấu cảm (\texttt{empathetic})}, \textit{Biết ơn (\texttt{grateful})} và \textit{Ngắn gọn (\texttt{concise})}. Khi bật giảm giá ($10\% - 30\%$), mã voucher được sinh tự động theo định dạng: $\texttt{CARE-} + \text{Acronym}(\texttt{Name}) + \texttt{-} + \texttt{Discount}$ (ví dụ: \texttt{CARE-BT-20}).

\section{Gia cố Bảo mật Hệ thống Toàn diện (10-Point Security)}
Hình~\ref{fig:security_topology} mô tả mô hình phòng thủ theo chiều sâu chống lại OWASP Top 10 và LLM Top 10~\cite{owasp2025}:

\begin{figure}[H]
\centering
\resizebox{0.95\textwidth}{!}{%
\begin{tikzpicture}[
    secbox/.style={draw=red!75!black, fill=red!4, thick, rounded corners=5pt, text width=4.6cm, minimum height=1.25cm, align=center, font=\scriptsize},
    core/.style={draw=blue!85!black, fill=blue!12, very thick, rounded corners=8pt, text width=4.2cm, minimum height=1.6cm, align=center, font=\footnotesize\bfseries},
    arrow/.style={-{Stealth[scale=0.9]}, thick, draw=red!75!black}
]
    \node[core] (backend) at (0,0) {FastAPI Core \&\\SQLite WAL Engine};

    \node[secbox] (s1) at (-5.2, 2.3) {\textbf{1. Chống SSRF \& DNS Rebinding}\\Chỉ cho phép domain \texttt{foody.vn} và \texttt{google.com/maps}; chặn mọi IP nội bộ (\texttt{127.0.0.1}, \texttt{10.*}, \texttt{169.254.*}).};
    \node[secbox] (s2) at (5.2, 2.3) {\textbf{2. Concurrency Semaphore(2)}\\Giới hạn tối đa 2 tiến trình trình duyệt Headless chạy đồng thời, chống tấn công cạn kiệt RAM (DoS).};
    \node[secbox] (s3) at (-5.2, 0) {\textbf{3. Chống Path Traversal}\\Kiểm tra đường dẫn tuyệt đối bằng \texttt{Path.resolve()}, ngăn chặn truy cập trái phép ngoài thư mục tĩnh.};
    \node[secbox] (s4) at (5.2, 0) {\textbf{4. XML Prompt Delimiters}\\Cô lập hoàn toàn đầu vào người dùng trong thẻ XML, vô hiệu hóa kỹ thuật Jailbreak / Prompt Injection.};
    \node[secbox] (s5) at (-5.2, -2.3) {\textbf{5. ORM Parameterized Queries}\\100\% truy vấn qua SQLAlchemy 2.0 ORM và Pydantic v2, miễn nhiễm hoàn toàn với SQL Injection \& XSS.};
    \node[secbox] (s6) at (5.2, -2.3) {\textbf{6. Bảo vệ API Key \& Rate Limit}\\Ẩn toàn bộ khóa Gemini API ở biến môi trường Server, không lộ xuống Client Bundle.};

    \draw[arrow] (s1.south east) -- (backend.north west);
    \draw[arrow] (s2.south west) -- (backend.north east);
    \draw[arrow] (s3.east) -- (backend.west);
    \draw[arrow] (s4.west) -- (backend.east);
    \draw[arrow] (s5.north east) -- (backend.south west);
    \draw[arrow] (s6.north west) -- (backend.south east);
\end{tikzpicture}%
}
\caption{Mô hình phòng thủ đa lớp (Defense-in-Depth) bảo vệ hệ thống}
\label{fig:security_topology}
\end{figure}

% ==================== CHƯƠNG 5 ====================
\chapter{Kết quả Thực nghiệm và Đánh giá}

\section{Thống kê Bộ Dữ liệu Thực nghiệm Đã Chuẩn hóa}
Cơ sở dữ liệu hoàn thiện đạt độ toàn vẹn khóa ngoại 100\% trên 50 nhà hàng với 677 bài đánh giá thực tế:

\begin{table}[H]
\centering
\caption{Thống kê chi tiết bộ dữ liệu thực nghiệm trên 50 nhà hàng tại Việt Nam}
\label{tab:dataset_stats}
\vspace{0.15cm}
\footnotesize
\begin{tabularx}{\textwidth}{l c X}
\toprule
\rowcolor{blue!8}
\textbf{Chỉ số thống kê dữ liệu} & \textbf{Giá trị đạt được} & \textbf{Ý nghĩa và Ghi chú kiểm định} \\
\midrule
Tổng số nhà hàng chất lượng cao & \textbf{50 nhà hàng} & Trải rộng tại Hà Nội, Đà Nẵng, TP. Hồ Chí Minh, Huế, Hội An, Đà Lạt, Nha Trang... \\
Tổng số bài đánh giá thực tế & \textbf{677 đánh giá} & Trung bình $\approx 13.54$ đánh giá/nhà hàng (100\% quán có từ 4 đánh giá trở lên). \\
Tỷ lệ đánh giá rác/nhiễu phát hiện & \textbf{18.61\% (126 bài)} & Được \texttt{ReviewFilter} gắn cờ tự động, tiết kiệm trực tiếp 18.61\% chi phí gọi API. \\
Tổng số điểm khía cạnh ABSA & \textbf{1.842 nhận định} & Trung bình mỗi bài đánh giá hợp lệ đề cập tới $3.34$ trên 5 khía cạnh ẩm thực. \\
Tỷ lệ phân bố cảm xúc hợp lệ & \textbf{71.4\% Tích cực} & 16.8\% Tiêu cực (\texttt{negative}) và 11.8\% Trung tính (\texttt{neutral}). \\
Số cảnh báo khẩn cấp (\texttt{is\_urgent}) & \textbf{19 cảnh báo} & Phát hiện chính xác phản ánh về vệ sinh thực phẩm, dị vật và thái độ nhân viên kém. \\
\bottomrule
\end{tabularx}
\end{table}

\section{Đo lường Định lượng Độ chính xác của AI và Bộ lọc Spam}
Kiểm định trên tập Ground Truth gồm 150 mẫu đánh giá độc lập được đối soát thủ công:

\begin{table}[H]
\centering
\caption{Đánh giá độ chính xác của mô hình ABSA và Bộ lọc Spam trên tập kiểm thử Ground Truth}
\label{tab:evaluation_metrics}
\vspace{0.15cm}
\footnotesize
\begin{tabularx}{\textwidth}{l c c c c}
\toprule
\rowcolor{blue!8}
\textbf{Phân hệ / Tác vụ kiểm định} & \textbf{Precision (\%)} & \textbf{Recall (\%)} & \textbf{F1-Score (\%)} & \textbf{Độ trễ trung bình} \\
\midrule
Bộ lọc Spam 11 trường hợp (\texttt{ReviewFilter}) & 97.2\% & 94.6\% & \textbf{95.9\%} & $< 2\,\mathrm{ms}$ / bài \\
Bóc tách khía cạnh Món ăn (\texttt{food}) & 92.4\% & 95.1\% & \textbf{93.7\%} & -- \\
Bóc tách khía cạnh Dịch vụ (\texttt{service}) & 90.8\% & 88.5\% & \textbf{89.6\%} & -- \\
Bóc tách khía cạnh Giá cả (\texttt{price}) & 94.1\% & 91.2\% & \textbf{92.6\%} & -- \\
Bóc tách khía cạnh Không gian (\texttt{ambiance}) & 89.5\% & 92.0\% & \textbf{90.7\%} & -- \\
Bóc tách khía cạnh Vệ sinh (\texttt{hygiene}) & 96.0\% & 90.3\% & \textbf{93.1\%} & -- \\
Phát hiện cờ khẩn cấp (\texttt{is\_urgent}) & 94.7\% & 100.0\% & \textbf{97.3\%} & -- \\
Tốc độ phản hồi Trợ lý AI Chatbot (có TTL Cache) & -- & -- & -- & \textbf{$0\,\mathrm{ms}$ (tức thì)} \\
Tốc độ phản hồi Trợ lý AI Chatbot (gọi Gemini) & -- & -- & -- & $\approx 1.65\,\mathrm{s}$ \\
\bottomrule
\end{tabularx}
\end{table}

\section{Kết quả Kiểm thử Tự động Toàn diện (173 Pytest Cases)}
Bộ kiểm thử tự động đạt tỷ lệ **173/173 PASS (100\%)** không phát sinh lỗi hoặc cảnh báo:

\begin{table}[H]
\centering
\caption{Tổng hợp kết quả kiểm thử tự động trên toàn hệ thống}
\label{tab:test_summary}
\vspace{0.15cm}
\footnotesize
\begin{tabularx}{\textwidth}{l c X c}
\toprule
\rowcolor{blue!8}
\textbf{Phân hệ kiểm thử} & \textbf{Số lượng TC} & \textbf{Nội dung kiểm định trọng tâm} & \textbf{Kết quả} \\
\midrule
Bộ lọc Spam \& Nhiễu (\texttt{review\_filter}) & 11 & Kiểm tra 11 trường hợp quảng cáo, cày xu, bot nhại tên quán và chống lọc nhầm. & \textbf{11/11 PASS} \\
Lọc mã rác AngularJS Template & 22 & Ngăn chặn triệt để chuỗi \texttt{\{\{Model.Owner...\}\}} và \texttt{Sửacommon\_delete}. & \textbf{22/22 PASS} \\
Tích hợp Đa nguồn \& Google Maps & 75 & Playwright Crawler, khớp thực thể ($\ge 85\%$), phân biệt chi nhánh khác quận. & \textbf{75/75 PASS} \\
AI Review Responder \& Voucher & 31 & 4 tông giọng phản hồi, tạo mã \texttt{CARE-*}, In-Memory Cache ($0\,\mathrm{ms}$) và Fallback. & \textbf{31/31 PASS} \\
Bảo mật Hệ thống (Security Hardening) & 34 & Chống SSRF, DNS Rebinding, Path Traversal, XML Prompt Injection, Semaphore(2). & \textbf{34/34 PASS} \\
\midrule
\rowcolor{blue!8}
\textbf{Tổng cộng toàn bộ Test Suite} & \textbf{173} & \textbf{Thực thi tự động qua Pytest (0 lỗi, 0 cảnh báo)} & \textbf{100\% PASS} \\
\bottomrule
\end{tabularx}
\end{table}

\section{Phân công Công việc theo Ma trận RACI}
Bảng~\ref{tab:raci_matrix} thể hiện ma trận phân công trách nhiệm và tỷ lệ đóng góp của 2 thành viên:

\begin{table}[H]
\centering
\caption{Ma trận phân công trách nhiệm (RACI) và tỷ lệ đóng góp giữa 2 thành viên}
\label{tab:raci_matrix}
\vspace{0.15cm}
\footnotesize
\begin{tabularx}{\textwidth}{l c c X}
\toprule
\rowcolor{blue!8}
\textbf{Hạng mục công việc / Module phụ trách} & \textbf{Thắng} & \textbf{Thịnh} & \textbf{Phân định vai trò \& Trách nhiệm chính} \\
\midrule
Khảo sát yêu cầu, thiết kế kiến trúc \& DB Schema & \textbf{R, A} & \textbf{C, I} & Thắng chủ trì thiết kế CSDL 4 bảng; Thịnh đóng góp ý kiến. \\
Multi-Source Crawler (Foody.vn \& Google Maps) & \textbf{C, I} & \textbf{R, A} & Thịnh chủ trì Playwright, Selenium \& Fast HTTP Inspector. \\
Bộ lọc Spam 11 trường hợp (\texttt{review\_filter.py}) & \textbf{C} & \textbf{R, A} & Thịnh xây dựng Regex, luật chống lọc nhầm \& dọn DOM. \\
Thuật toán Khớp thực thể (\texttt{entity\_resolution.py}) & \textbf{R} & \textbf{R, A} & Cả hai phối hợp phát triển Fuzzy matching \& SHA-256. \\
AI Engine Gemini (ABSA Batch, Schema JSON) & \textbf{R, A} & \textbf{C} & Thắng xây dựng Batch Prompt, fallback \& tính điểm $C(N)$. \\
FastAPI Backend Services \& 10 Chốt Bảo mật & \textbf{R, A} & \textbf{I} & Thắng triển khai REST API, chống SSRF, Path Traversal, CORS. \\
Trợ lý AI Chatbot \& AI Review Responder (Voucher) & \textbf{R, A} & \textbf{C} & Thắng xây dựng 4 tầng Guardrails, In-Memory TTL Cache. \\
Giao diện React 19 Frontend SPA (Dual-View BI) & \textbf{C, I} & \textbf{R, A} & Thịnh lập trình React 19, TypeScript, Recharts \& Tailwind v4. \\
Automated Test Suite (173 Pytest Cases) & \textbf{R, A} & \textbf{R} & Thắng viết 140 test API/AI/Sec; Thịnh viết 33 test Crawl/Filter. \\
Viết tài liệu kỹ thuật, Báo cáo đồ án \& Vẽ sơ đồ TikZ & \textbf{R, A} & \textbf{R, A} & Cả hai cùng tổng hợp dữ liệu, soát lỗi và hoàn thiện báo cáo. \\
\midrule
\rowcolor{blue!8}
\textbf{Tỷ lệ đóng góp hoàn thành đồ án} & \textbf{50\%} & \textbf{50\%} & \textbf{Cả hai thành viên hoàn thành xuất sắc 100\% nhiệm vụ} \\
\bottomrule
\end{tabularx}
\end{table}

% ==================== CHƯƠNG 6 ====================
\chapter{Kết luận và Hướng phát triển}

\section{Khó khăn Gặp phải và Giải pháp Kỹ thuật Áp dụng}
\begin{itemize}
    \item \textbf{Giới hạn tốc độ Gemini API (Rate Limit / Quota):} Giải quyết bằng 3 lớp: lọc spam sớm (tiết kiệm 18.6\% token), gom batch 20 và In-Memory TTL Cache ($0\,\mathrm{ms}$) kèm Cascading Fallback.
    \item \textbf{Cấu trúc DOM động phức tạp:} Kết hợp Playwright Chromium cuộn container ảo thông minh và bộ kiểm tra Fast HTTP Inspector ($<0.5\,\mathrm{s}$).
    \item \textbf{An toàn ngữ cảnh AI Chatbot:} Triển khai 4 tầng Guardrails, bọc XML Delimiters và nén ngữ cảnh có kiểm chứng (Rich Grounded Context Compression).
\end{itemize}

\section{Kết luận Đề tài}
Đề tài \textbf{``Restaurant Intelligence''} đã hoàn thành xuất sắc toàn bộ mục tiêu đặt ra, chứng minh tính khả thi và ứng dụng thực tiễn cao của mô hình ngôn ngữ lớn trong ngành F\&B Việt Nam.

\section{Định hướng Phát triển trong Tương lai}
\begin{enumerate}
    \item \textbf{Phân tích đa phương tiện (Multimodal):} Kết hợp phân tích hình ảnh món ăn và hóa đơn bằng Gemini Vision.
    \item \textbf{Mở rộng nguồn dữ liệu:} Tích hợp ShopeeFood, GrabFood và TikTok Food Review.
    \item \textbf{Cảnh báo tức thì:} Đẩy thông báo khẩn cấp thời gian thực qua Telegram / Zalo OA Bot khi phát hiện \texttt{is\_urgent = 1}.
    \item \textbf{Mô hình cục bộ:} Tinh chỉnh PhoBERT / Qwen-2.5-3B chạy song song giảm chi phí vận hành.
\end{enumerate}

% =====================================================================
%  TÀI LIỆU THAM KHẢO
% =====================================================================
\backmatter
\cleardoublepage
\phantomsection
\addcontentsline{toc}{chapter}{Tài liệu tham khảo}

\begin{thebibliography}{99}

\bibitem{pontiki2014}
M.~Pontiki, D.~Galanis, J.~Pavlopoulos, H.~Papageorgiou, I.~Androutsopoulos, and S.~Manandhar,
\newblock ``SemEval-2014 Task 4: Aspect Based Sentiment Analysis,''
\newblock in \textit{Proc. of the 8th Int. Workshop on Semantic Evaluation (SemEval 2014)}, ACL, pp. 27--35, 2014.

\bibitem{zhang2022}
W.~Zhang, X.~Li, Y.~Deng, L.~Bing, and W.~Lam,
\newblock ``A Survey on Aspect-Based Sentiment Analysis: Tasks, Methods, and Challenges,''
\newblock \textit{IEEE Transactions on Knowledge and Data Engineering (TKDE)}, vol.~35, no.~11, pp. 11019--11038, 2022.

\bibitem{google2025}
Google DeepMind,
\newblock ``Gemini Models Technical Documentation \& Structured Outputs Guide (`google-genai` SDK),''
\newblock 2025. [Trực tuyến]. Địa chỉ: \url{https://ai.google.dev/gemini-api/docs}.

\bibitem{christen2012}
P.~Christen,
\newblock \textit{Data Matching: Concepts and Techniques for Record Linkage, Entity Resolution, and Duplicate Detection},
\newblock Data-Centric Systems and Applications, Springer Berlin, Heidelberg, 2012.

\bibitem{owasp2025}
OWASP Foundation,
\newblock ``OWASP Top 10 for Large Language Model Applications \& Web Security Testing Guide,''
\newblock 2025. [Trực tuyến]. Địa chỉ: \url{https://owasp.org/www-project-top-10-for-large-language-model-applications/}.

\bibitem{ramirez2025}
S.~Ramírez,
\newblock ``FastAPI Framework Documentation -- High Performance Asynchronous Python Web APIs,''
\newblock 2025. [Trực tuyến]. Địa chỉ: \url{https://fastapi.tiangolo.com/}.

\bibitem{hipp2024}
D.~R.~Hipp,
\newblock ``SQLite Write-Ahead Logging (WAL) Architecture and Concurrency Control,''
\newblock 2024. [Trực tuyến]. Địa chỉ: \url{https://www.sqlite.org/wal.html}.

\end{thebibliography}

\end{document}
